\section{Métodos}
\label{sec:method}

\paragraph{Dados e splits.}
O conjunto tem 4946 imagens de cinco categorias e de background. O val fornecido é o nosso test. Do train fornecido separámos 20\,\% para validação, por imagem (há imagens com várias categorias), com estratificação pela categoria mais rara e seed~42 (Tab.~\ref{tab:splits}). Cada categoria é uma
tarefa binária, com os seus positivos contra o background. As outras categorias ficam fora do treino e do ranking. Esta escolha segue a Stage~B e não é one-vs-rest. Todos os métodos usam as mesmas listas e labels.

\begin{table}[t]
\centering
\caption{Imagens por categoria e split. A última linha é a prevalência no
test, ou seja, o AP de um ranking aleatório (mAP 0.209).}
\label{tab:splits}
\small
\setlength\tabcolsep{4pt}
\begin{tabular}{lrrrrrr}
\toprule
 & aero & mbike & person & horse & car & backgr. \\
\midrule
train & 90  & 96  & 821 & 111 & 298 & 815  \\
val   & 22  & 24  & 204 & 28  & 78  & 204  \\
test  & 126 & 125 & 983 & 148 & 337 & 1077 \\
\midrule
prev.\ & .105 & .104 & .477 & .121 & .238 & -- \\
\bottomrule
\end{tabular}
\end{table}

\paragraph{Features handcrafted.}
As imagens passam a cinzento com altura máxima de 480\,px. O dense SIFT (\texttt{vl\_phow}, passo 4\,px, bins 4--10\,px) é reduzido a 80-d por PCA para VLAD e FV. Os vocabulários e os GMM (covariância diagonal) são ajustados só com imagens de treino (733 imagens, até 400 descritores cada). Cada imagem
é codificada em tiles $2\times2$ concatenados. A variante global soma os tiles, o que só é exato no BoVW. Comparamos nenhuma normalização, L1, L2 e Hellinger no BoVW, e nenhuma, L2, \emph{signed square root} e intra-normalização (só VLAD) na VLAD e no FV.

\paragraph{Features profundas.}
Usamos o pré-processamento oficial de cada modelo: resize para 256 e crop de 224 na ResNet-18, no ViT-B/32 (torchvision \texttt{IMAGENET1K\_V1}) e no DINOv2 ViT-S/14, e o \texttt{preprocess} próprio do CLIP ViT-B/32 (224, bicúbico). As features são o pooling de \texttt{layer2}--\texttt{4} da ResNet, o token de classe do ViT, o CLS ou a média dos patches do DINOv2 e o embedding de imagem do CLIP, todos normalizados em L2 (escolhido na validação). O few-shot e o estudo de $C$ usam o CLS do DINOv2.

\paragraph{Classificadores e seleção.}
Usamos um \texttt{LinearSVC} por categoria ($C=10$, $\texttt{max\_iter}=100$, como no código fornecido). Com este limite, o BoVW-Hellinger, o FV em aeroplane e quase todos os SVMs sobre features profundas não convergem. Nestas últimas, convergir (até 20\,000 iterações) muda o mAP de test no máximo 0.0006. O CLIP zero-shot ordena cada categoria pela similaridade de cosseno ao seu prompt. As escolhas são feitas pelo mAP de validação, e o test é avaliado uma vez. As exceções são a Parte~3, que reporta todas as condições, e a Stage~D da Parte~4, que escolhe $C$ para 5 positivos e para todos os positivos.

\paragraph{Métricas.}
O AP é o do scikit-learn (sem interpolação), calculado sobre o score contínuo. O mAP é a média das 5 categorias. Reportamos também precisão e recall no threshold~0 e, na Parte~5, P@K e R@K. Nas matrizes de confusão vence o maior score (background se todos forem negativos; no zero-shot, o argmax entre 6 prompts), e excluímos as imagens com várias categorias (restam 2169 no test).

\paragraph{Repetições, hardware e custos.}
Reportamos média $\pm$ desvio padrão com 5 seeds na Stage~G, 3 na Parte~3 e 5 sorteios no few-shot. Tudo correu num Ryzen~5 5500, sem GPU, em Windows (Python 3.11.9, scikit-learn 1.9.1, PyTorch 2.14.0, cyvlfeat 0.7.1). A extração profunda inclui o carregamento do modelo e o pré-processamento. A handcrafted foi medida em 5 imagens, por isso só se compara em ordem de
grandeza. O armazenamento é dimensão $\times$ 4\,B. Os scripts estão em \texttt{runs/}.